\section{Discussion}
\label{sec:discussion}

% The numerical results answer the motivating question conditionally: short
% pulse sequences can contribute directly to frequency calibration, but they do
% not automatically become general-purpose frequency estimators.  The present
% two-branch sequence is a local probe, not a wide-range replacement for Ramsey
% spectroscopy.  Its cubic response reduces local nonlinear bias but cannot
% remove a response fold.  In closed loop, the useful range is preserved by
% changing the measurement reference so that the short-pulse probe stays local;
% merely assigning it to a coarse stage before an expensive fine estimator does
% not guarantee a lower total cost.

% This framing clarifies what it means for a pulse sequence to participate in
% calibration.  The sequence must encode the parameter in an observable response,
% that response must admit a calibrated inverse over a stated operating window,
% and the surrounding protocol must keep subsequent measurements inside that
% window.  In the present example, the qubit experiences constant detuning
% throughout two finite pulses, so the third-order coefficient contains
% unequal-time contributions.  A time-diagonal kernel is not the relevant
% object.  An odd-polynomial scan can recover the same coefficient when centered
% at the symmetry point, but that assumption is fragile under a moving drive.
% Evaluating the full response kernel at the actual reference frequency makes
% the short-pulse estimator and controller definitions consistent.

% The cost comparison is intentionally based on instrumented solver calls rather
% than a nominal workflow price.  The nominal price counts the two population
% measurements but can omit the solver work required to recalibrate a response
% coefficient when the drive changes.  Counting every solver invocation avoids a
% large artificial speedup.  Wall-clock time supports the same ordering in the
% present implementation, but it depends on hardware, caching, and solver choice
% and is not a universal physical metric.

% Several steps are needed before interpreting the protocol as an experimental
% calibration advantage.  The current simulations use ideal control envelopes,
% projective readout, a Markovian two-level model, and deterministic expectation
% values. Their reported uncertainty is therefore zero by construction, not a
% measured confidence interval. Shot noise, assignment error, drift during acquisition, pulse
% distortion, leakage, and imperfect knowledge of the kernel will perturb both
% the detuning estimate and the secant slope.  An experiment should report the
% number of circuit repetitions and elapsed acquisition time, and it should
% recompute the cost-to-reach curve at a common residual threshold.  A wide-range
% capture or range-guard path is also required when no reliable initial frequency
% seed is available.

% These limitations suggest a practical hierarchy rather than a replacement
% claim.  Spectroscopy or a Ramsey scan supplies acquisition and recovery.  Once
% the transition frequency is known, short-pulse sequences can participate in
% frequent local maintenance, here through a drive-tracked two-branch
% discriminator.  Occasional wide-range measurements can reinitialize the local
% loop after a range violation.  This division uses pulse-time information where
% it is unambiguous while retaining established estimators for global capture.

% The implemented state machine makes this hierarchy operational by separating
% \textsc{Verify} from \textsc{Lock}. It freezes the candidate bias during
% independent verification, sends monitor suspicion and periodic Ramsey audits
% from \textsc{Lock} through \textsc{Verify}, and reserves direct reacquisition
% for large jumps or loss of reference. These are protocol and software
% capabilities, not evidence that long-term stability has been achieved. No
% hardware time series is presently available from which to infer closed-loop
% bandwidth, frequency-noise suppression, a power spectral density, or Allan
% stability.
